\section{A sharper finite natural-dimension bound for Spin cores}
\label{sec:spin-sharpened}

\paragraph{Status and scope.}
This is the separate new quantitative derivation D-OL-25--26.
Independent whole-proof reviews are pending.  It preserves the
frozen mathematical texts D-OL-19--24 and uses their complete
Clifford, radical, normalizer and actual-source arguments.
All fields here have odd characteristic.  Modules may have
arbitrary central kernels; no descent to an ordinary natural
orthogonal group is assumed.

Put
\[
 \beta_0=\frac16-\frac1{10^7}
         =\frac{4999997}{30000000},\qquad
 \delta_0=\frac16+\frac1{10^5}
         =\frac{50003}{300000},\qquad r_0=\frac{419}{53}.
 \tag{SSP.1}
\]

\subsection*{D-OL-25. The exact support threshold and the sharper Gauss bound}

\begin{lemma}[D-OL-25]
Let $q$ be odd, and let $M$ be any nontrivial irreducible complex
module of $S=\operatorname{Spin}_d(q)$, of degree $a$.
Suppose $d\ge751$ is odd or $d\ge752$ is even.
For every nonscalar compatible finite-projective-order normalizer $T$,
\[
 \nu(T)\ge\beta_0a,
 \qquad q\ge5\ \Longrightarrow\ \nu(T)>\delta_0a.
 \tag{SSP.2}
\]
If $q=3$ and $\nu(T)\le\delta_0a$, the natural projective
lift is a reflection and has projective order two.
In odd dimension its unique $SO$ representative is a negative
reflection.  Moreover $a>10^6$.
\end{lemma}

\begin{proof}
We first check the published character input at the new threshold.
In the notation of Larsen--Shalev--Tiep, formula (4.3.2) has
the following \emph{second} bound:
\[
 |\chi(s)|/\chi(1)<(\gamma_q)^{\sqrt{N/6}}
        \quad\hbox{if }\operatorname{supp}_{\mu_2}(s)\ge N\ge144.
 \tag{SSP.3}
\]
Here $s\in\operatorname{Spin}_d(q)$ and $\chi$ is any nontrivial
irreducible character; the support is that of its natural image.
The restricted scalar set $\mu_2$ is permitted by Remark~4.3.7.
The exact published locators are
\cite[(4.1.1), p.~1907; Theorem~4.3.6, (4.3.2), p.~1918;
complete proof pp.~1918--1923; Remark~4.3.7, p.~1923]{HOS:LST2011}.
Thus (SSP.3) is not a substitution into the different first
line $\sqrt{N/7}$, whose displayed threshold is $567$.
For every $q$ under consideration, $\gamma_q\le19/20$:
it equals $19/20$ for $q<43$, and is
$1/(\sqrt q-1)<1/5$ for $q\ge43$.

At $N=375$, all constants can be checked by integer arithmetic.
In fact
\[
 r_0^2=\frac{175561}{2809}<\frac{125}{2}=\frac{375}{6},
 \qquad 2\cdot175561=351122<351125=125\cdot2809,
 \tag{SSP.4}
\]
and the exact positive-integer comparison is
\[
                  19^{419}150000^{53}
                    <20^{419}99997^{53}.                    \tag{SSP.5}
\]
The supplied certificate checks (SSP.5) directly by integer
exponentiation; the two integers have $2692$ bits.
Taking positive $53$rd roots and using that $0<19/20<1$ gives
\[
 (19/20)^{\sqrt{375/6}}
       <(19/20)^{419/53}
       <\frac{99997}{150000}=1-2\delta_0.
 \tag{SSP.6}
\]
There is no floating-point estimate in this implication.
For a finite-order complex operator, a largest-eigenspace
eigenvalue can be rotated to one, giving
$|\operatorname{tr}T|\ge a-2\nu(T)$.
Consequently (SSP.3)--(SSP.6) prove
\[
 \operatorname{supp}_{\mu_2}(s)\ge375
            \quad\Longrightarrow\quad\nu(\rho(s))>\delta_0a.
 \tag{SSP.7}
\]
A central choice of Spin lift changes no natural support
or absolute character value.

We next sharpen only the rank-one radical estimate.
Let $W$ be a nondegenerate quadratic space of dimension $m\ge43$
over $\mathbf F_q$, and let a reflection act on the dual of
the abelian full-line radical.  A nonzero Levi orbit is a
nonzero norm fiber in $W$.  If its prescribed norm is $c$,
the elementary quadratic Gauss count gives
\[
 |\{w\in W:Q(w)=c\}|\ge q^{m-1}-q^{m/2}.
 \tag{SSP.8}
\]
Indeed the Fourier expansion over the $q$ additive characters
has trivial term $q^{m-1}$ and its other $q-1$ terms each
have absolute value $q^{m/2-1}$.  The estimate applies
to every $c$, including zero.  Removing a possible zero vector
decreases the lower bound by at most one.
The fixed hyperplane is nondegenerate because the reflecting
normal is nonsingular.  The same count in dimension $m-1$
bounds its prescribed-norm vectors above by
$q^{m-2}+q^{(m-1)/2}$.
Thus the fraction of fixed nonzero weights in any such orbit is at most
\[
 \frac{q^{m-2}+q^{(m-1)/2}}
      {q^{m-1}-q^{m/2}-1}
 \le c_q^*:=\frac{q^{-1}+q^{-20}}
                      {1-q^{-20}-q^{-40}}.
 \tag{SSP.9}
\]
For the second inequality divide numerator and denominator
by $q^{m-1}$.  Since $m\ge43$, the numerator's error
$q^{-(m-1)/2}$ is at most $q^{-21}\le q^{-20}$;
the denominator's two errors $q^{1-m/2}$ and $q^{1-m}$
are at most $q^{-20}$ and $q^{-40}$, respectively.
All denominators are positive.

The right side decreases with $q$: its positive numerator
decreases and its positive denominator increases.
The exact endpoint comparisons are
\[
 c_3^*<\frac13+\frac4{10^7}=1-4\beta_0,
 \qquad c_5^*<\frac14.
 \tag{SSP.10}
\]
They are rational inequalities checked by clearing their
displayed positive integer denominators in the certificate.
The opposite lifted-radical fixed space has dimension at
most $a/2$.  All nonfixed weight orbits of the involution
have length two.  Hence the support fraction is at least
$(1-c_q^*)/4$, which is greater than $\beta_0$ for $q=3$
and greater than $3/16>\delta_0$ for $q\ge5$.
The full inverse-image Levi supplies the same orbit
multiplicities for every represented Spin quotient, as proved
in D-OL-19--20.  No natural-group extension of the module
has been used in (SSP.8)--(SSP.10).

For completeness, we identify every use of the old natural
dimension threshold in the full normalizer proof.
The Clifford lift and exact central-kernel compatibility
need only $d\ge22$.
The elementary fixed-hyperbolic-pair, quadratic-field and
semilinear descent cases of D-OL-20 and D-OL-23 give support
at least $3a/16$ or $7a/32$, apart from the rank-one case
just sharpened.  Their radical complements have dimension
at least $d-4$, so their hypotheses are still satisfied.
In a prime-order case without a fixed hyperbolic pair,
the same polar-rank argument gives restricted natural
support at least the integer ceiling of $(d-2)/2$.
This is at least $375$ at both new thresholds.
Equation (SSP.7) therefore applies in exactly the place
where the earlier proof used support $567$.
No smaller-dimensional restriction or ordinary-character
degree bound is inserted.

The all-orders reductions are preserved as well.
An odd prime projective power is covered by the preceding
cases.  Over $\mathbf F_3$ there is no field component.
In even natural dimension, the complete determinant and
Clifford compatibility argument of D-OL-20 excludes
non-reflection low-support orders; its only character-estimate
invocations are those just strengthened.
In odd natural dimension, an order $2^s$ lift with $s\ge3$
has its square in $\Omega$, with both $1$ and $-1$
eigenspaces contained in the one-dimensional reflecting
line, as in (OSP.18).  Its restricted support is at least
$d-1>375$, so (SSP.7) applies to a Spin lift of the square.
The remaining fourth-order case is D-OL-22, which needs
only $d\ge25$ and gives $7a/32>\delta_0a$.
Thus the only possible low-support operators are precisely
the projective reflections stated in the lemma.

Finally the all-Spin degree argument in D-OL-20 and D-OL-23
gives, with $m=d-2\ge749$,
\[
 \begin{split}
 a&\ge q^{m-1}-q^{\lceil m/2\rceil}-1\\
  &\ge(q-1)q^{m-2}-1
   \ge2\cdot3^{747}-1>10^6.
 \end{split}                                                \tag{SSP.11}
\]
Here $\lceil m/2\rceil\le m-2$; the last strict inequality
already holds after replacing the exponent $747$ by $15$.
The degree argument uses a nontrivial lifted-radical weight
and its full-preimage Levi orbit.  It therefore covers genuine
Spin characters with every central kernel.
This completes all assertions.
\end{proof}

\subsection*{D-OL-26. The actual mixed-core source consequence}

\begin{theorem}[D-OL-26]
An irreducibly acting nonempty normal central product of
represented odd-characteristic Spin factors cannot occur in
a rank-nine six-dimensional multiplicity source if every odd
natural dimension is at least $751$ and every even natural
dimension is at least $752$.
All irreducible factor modules, represented central quotients,
compatible scalars, outer automorphisms and factor permutations
are allowed.  No transitivity or selfduality is required.
In particular, in the transitive-core source problem the
remaining odd natural dimension is at most $749$ and the
remaining even natural dimension is at most $750$.
\end{theorem}

\begin{proof}
Let $e$ be the product of the irreducible factor degrees.
By (SSP.11), $e>10^6$.
Retain the complete good/bad and generation proof of
D-OL-21 and D-OL-24, now defining good to mean actual
moving rank at least $\delta_0e$.
Every genuine inertia moves at least $\beta_0e$.
The weighted factor-permutation bound is at least $e/4$,
which exceeds $\delta_0e$, and nonidentity scalars move all $e$.
For a factor-preserving bad inertia, each active factor is
a projective reflection over $\mathbf F_3$ by D-OL-25.
The two normalized eigenspace proportions of an active
factor lie in $[\beta_0,1-\beta_0]$.
Two active factors have support at least
$2\beta_0(1-\beta_0)e>e/4>\delta_0e$.
Thus every bad inertia has exactly one active factor.
This retains arbitrary original scalar phases and arbitrary
represented Spin centers.

The Euler identity and actual finite-cover moduli argument
are unchanged: a proposed rank-nine source must have a
sphere intermediate base, at least twelve genuine marks,
a generating product-one inertia tuple and total moving
rank $D=2e+6$.
The exact numerical margins are
\[
 \begin{gathered}
 13\beta_0-2>\frac17,\qquad
 11\beta_0+\delta_0-2=\frac{89}{10^7},\\
 2\beta_0(1-\beta_0)>\frac14,\qquad
 \frac{6\cdot10^7}{89}<10^6.
 \end{gathered}                                             \tag{SSP.12}
\]
Thirteen marks would give $e<42$; twelve marks including
one good inertia would give $e\le6\cdot10^7/89<10^6$.
Both contradict (SSP.11).  Exactly twelve bad inertias
therefore remain.

All preserve the individual factors.  A factor over a
field different from $\mathbf F_3$ has trivial inertia
projection, contradicting generation of its inner simple
group.  In a factor over $\mathbf F_3$, at most twelve
projective reflections would generate a projective group
containing the full natural $P\Omega$.
Choose natural reflection representatives.  Their common
fixed subspace has dimension at least $d-12>0$.
Every natural $\Omega$ element is projectively a word
in those reflections, so it acts by a scalar on that
subspace.  Perfectness makes this scalar action trivial,
contradicting the absence of a nonzero natural fixed
space for $\Omega$.  This is the same actual-generation
argument in either parity; in even dimension it is made
after the natural projection to $P\Omega$, so no splitting
of the natural or represented Spin center is asserted.
The contradiction proves the theorem.
\end{proof}

\paragraph{Certificate and exact boundary.}
The accompanying finite certificate checks every rational
margin, the positive-integer comparison (SSP.5), the minimal
support dimensions and the degree floor.  The displayed
Gauss-sum and monotonicity arguments prove the all-dimension
claims; no finite parameter scan replaces them.
This is a strict reduction of the necessary parameter range,
not an existence claim at any remaining dimension.
